\subsection{有理直线}

考察集 \(\mathfrak{F}\) ：它由对称开区间 \(]-a, +a[\) （其中 \(a\) 取遍 \(>0\) 的有理数组成的集）组成；我们将证明 \(\mathfrak{F}\) 是 \(o\) 在一个与 \(\mathbf{Q}\) 的加法群结构相容的拓扑中的基本邻域系。

群 Q 是交换的，且公理 \((\mathbf{GV}_{\mathrm{II}}^{\prime})\) 显然满足；因此只需证明公理 \((\mathbf{GV}_{\mathrm{I}}^{\prime})\) 也满足，换句话说，对每个 a \textgreater{} 0 ，存在 b \textgreater{} 0 ，使得条件 \(|x| < b\) 与 \(|y| < b\) 合起来蕴含 \(|x + y| < a\) 。三角不等式表明可取 b = a/2。

定义 1. 有理直线是指由集合 Q 连同如下加法群拓扑所构成的拓扑空间：对称开区间 \(]-a, +a[\) （a \textgreater{} o）组成 o 的一个基本邻域系。

如此定义的拓扑群 Q 称为有理直线的加法群。

若 \(a\) 是任一 \(> 0\) 的有理数，则存在整数 \(n > 0\) 使得 \(1 / n < a\) ；因而开区间 \(] - \frac{1}{n}, + \frac{1}{n} [ (n = 1, 2, \ldots) \text{ 共同}\)

组成有理直线上 o 的一个基本邻域系。

对任一点 \(x \in Q\) ，取开区间 \(]x - a, x + a[\) （其中 a 取遍 \textgreater0 的有理数组成的集，或数 \(1/n\) 组成的集），就得到该点的一个基本邻域系。

因此定义 1 与第一章 § 1 第 2 目中给出的定义等价。

对每对有理数 \((a, b)\) （满足 a \textless{} b），存在 \(c \in Q\) 使得 a \textless{} c \textless{} b {[}例如 \(c = (a + b)/2\) {]}；由此可知有理直线是一个非离散的豪斯多夫空间。

对每个 \(a > 0\) ，令 \(\mathbf{U}_a\) 为由所有对 \((x, y)\) （属于 \(\mathbf{Q} \times \mathbf{Q}\) ，满足 \(|x - y| < a\) ）组成的集。当 \(a\) 取遍 \(> 0\) 的有理数组成的集（或只取数 \(1/n\) 组成的集）时，诸集 \(\mathbf{U}_a\) 组成有理直线加法群 \(\mathbf{Q}\) 的一致结构的一致邻域基本系。关系 (2) 与 (3) 表明 \(|x|, x^+\) 与 \(x^-\) 在 \(\mathbf{Q}\) 上一致连续。由此推出函数 \(\sup (x, y)\) 与 \(\inf (x, y)\) 在 \(\mathbf{Q} \times \mathbf{Q}\) 上一致连续。

